% =============================================================================
%  导言区：中文译本《局部同调镜像对称》
%  原文：arXiv:2611.07202v1 [math.SG], 4 Nov 2026, 56 pp.
%        Local Homological Mirror Symmetry
%        Si-Yang Liu (Imperial College London), Kyungmin Rho (MPIM Bonn)
%  编译：xelatex main.tex（两遍）
% =============================================================================
\documentclass[11pt,a4paper]{article}

\usepackage[UTF8, fontset=fandol, heading=true]{ctex}
\usepackage{amsmath}
\usepackage{amssymb}
\usepackage{amsthm}
\usepackage{mathtools}
\usepackage{bm}
\usepackage[a4paper,left=2.5cm,right=2.5cm,top=2.6cm,bottom=2.6cm]{geometry}
\usepackage{graphicx}
\usepackage{xcolor}
\usepackage[most]{tcolorbox}
\usepackage{enumitem}
\usepackage{booktabs}
\usepackage{longtable}
\usepackage{array}
\usepackage{tikz}
\usepackage{tikz-cd}
\usetikzlibrary{shapes.geometric,arrows.meta,positioning,calc}
\usepackage[font=small,labelfont=bf,labelsep=quad]{caption}
\usepackage{cite}
\usepackage[colorlinks=true,
            linkcolor=blue!45!black,
            citecolor=blue!45!black,
            urlcolor=blue!45!black]{hyperref}

\linespread{1.15}
\setlength{\emergencystretch}{3em}
\allowdisplaybreaks[3]
\pagestyle{plain}
\renewcommand{\figurename}{图}
\renewcommand{\tablename}{表}
\renewcommand{\refname}{参考文献}
\renewcommand{\contentsname}{目录}
\renewcommand{\appendixname}{附录}
\renewcommand{\proofname}{证明}
\renewcommand{\abstractname}{摘要}

% ---------------------------------------------------------------------------
% 原文字体约定：原稿把 \emph 重定义为无衬线粗体；中译本里改作粗体（中文加粗，
% 英文粗体），用于标记作者强调的关键术语。
% ---------------------------------------------------------------------------
\renewcommand{\emph}[1]{{\bfseries #1}}

% ---------------------------------------------------------------------------
% 原文数学宏（与原文 preamble 逐一对应，保证公式可逐符号对照）
% ---------------------------------------------------------------------------
\DeclareMathOperator*{\mref}{\mathsf{ref}}
\DeclareMathOperator*{\rref}{\mathsf{rref}}
\DeclareMathOperator{\Tr}{\mathsf{Tr}}
\DeclareMathOperator{\espan}{\mathrm{span}}
\DeclareMathOperator{\row}{\mathsf{row}}
\DeclareMathOperator{\col}{\mathsf{col}}
\DeclareMathOperator{\rank}{\mathrm{rank}}
\DeclareMathOperator{\Tot}{Tot}
\DeclareMathOperator{\Rng}{\mathsf{Rng}}
\DeclareMathOperator{\id}{\mathsf{Id}}
\DeclareMathOperator{\diag}{\mathsf{diag}}
\DeclareMathOperator{\Fun}{\mathsf{Fun}}
\DeclareMathOperator{\diff}{diff}
\DeclareMathOperator{\Gr}{\mathrm{Gr}}
\DeclareMathOperator{\QCoh}{\mathrm{QCoh}}
\DeclareMathOperator{\Spec}{\mathrm{Spec}}
\DeclareMathOperator{\Max}{\mathsf{Max}}
\DeclareMathOperator{\Mod}{\mathrm{Mod}}
\DeclareMathOperator{\End}{\mathrm{End}}
\DeclareMathOperator{\Nat}{\mathsf{Nat}}
\DeclareMathOperator*{\colim}{\mathrm{colim}}
\DeclareMathOperator*{\hocolim}{hocolim}
\DeclareMathOperator*{\holim}{holim}
\DeclareMathOperator{\Ab}{\mathrm{Ab}}
\DeclareMathOperator{\cl}{\mathsf{cl}}
\DeclareMathOperator{\tr}{\mathsf{tr}}
\DeclareMathOperator{\prim}{\mathrm{prim}}
\DeclareMathOperator{\codim}{\mathrm{codim}}
\DeclareMathOperator{\Res}{\mathrm{Res}}
\DeclareMathOperator{\Sets}{\mathsf{Sets}}
\DeclareMathOperator{\GL}{\mathrm{GL}}
\DeclareMathOperator{\Gal}{\mathrm{Gal}}
\DeclareMathOperator{\hs}{\mathfrak{hs}}
\DeclareMathOperator{\Cone}{\mathrm{Cone}}
\DeclareMathOperator{\Ext}{\mathrm{Ext}}
\DeclareMathOperator{\MHS}{\mathrm{MHS}}
\DeclareMathOperator{\rat}{\mathrm{rat}}
\DeclareMathOperator{\MHM}{\mathsf{MHM}}
\DeclareMathOperator{\Perv}{\mathsf{Perv}}
\DeclareMathOperator{\DMod}{\mathrm{DMod}}
\DeclareMathOperator{\Coh}{\mathsf{Coh}}
\DeclareMathOperator{\DR}{\mathrm{DR}}
\DeclareMathOperator*{\pt}{\mathrm{pt}}
\DeclareMathOperator{\Hdg}{\mathrm{Hdg}}
\DeclareMathOperator{\supp}{\mathrm{supp}}
\DeclareMathOperator*{\can}{\mathrm{can}}
\DeclareMathOperator{\var}{\mathrm{var}}
\DeclareMathOperator{\ord}{\mathrm{ord}}
\DeclareMathOperator{\Fuk}{\mathsf{Fuk}}
\DeclareMathOperator{\Perf}{\mathsf{Perf}}
\DeclareMathOperator{\IndCoh}{IndCoh}
\DeclareMathOperator{\Ch}{\mathsf{Ch}}
\DeclareMathOperator{\Aut}{Aut}
\DeclareMathOperator{\Obj}{Obj}
\DeclareMathOperator{\Ob}{Ob}
\DeclareMathOperator{\field}{\mathbb{k}}
\DeclareMathOperator{\Hom}{\operatorname{Hom}}
\DeclareMathOperator{\HHom}{\mathcal{H}om}
\DeclareMathOperator{\Proj}{\operatorname{Proj}}
\DeclareMathOperator{\Sym}{\operatorname{Sym}}
\DeclareMathOperator{\dg}{\operatorname{dg}}
\DeclareMathOperator{\Sch}{\mathsf{Sch}}
\DeclareMathOperator{\VB}{\mathsf{VB}}
\DeclareMathOperator{\vb}{\mathsf{vb}}
\DeclareMathOperator{\str}{\operatorname{str}}
\DeclareMathOperator{\weight}{\operatorname{weight}}
\DeclareMathOperator{\pr}{\operatorname{pr}}
\DeclareMathOperator{\Gm}{\mathbb{G}_{\operatorname{m}}}
\DeclareMathOperator{\cha}{\lambda}
\DeclareMathOperator{\chaa}{\mu}
\DeclareMathOperator{\comp}{\vcenter{\hbox{\scalebox{0.65}{$\bullet$}}}}
\DeclareMathOperator{\Nbd}{Nbd}
\DeclareMathOperator{\Bl}{Bl}
\DeclareMathOperator{\Symp}{Symp}
\DeclareMathOperator{\rF}{\reflectbox{F}}
\DeclareMathOperator{\objectdegree}{{\text{$i$}}}
\DeclareMathOperator{\homdegree}{{\text{$k$}}}
\DeclareMathOperator{\yo}{\text{よ}}
\DeclareMathOperator{\Rep}{Rep}
\DeclareMathOperator{\Ku}{\textrm{Kü}}
\DeclareMathOperator{\Sh}{Sh}
\newcommand{\W}{\mathcal{W}}
\newcommand{\Z}{\mathbb{Z}}
\newcommand{\R}{\mathbb{R}}
\newcommand{\C}{\mathbb{C}}
\newcommand{\N}{\mathbb{N}}
\newcommand{\Q}{\mathbb{Q}}
\newcommand{\kk}{\mathbb{k}}
\newcommand{\pp}{\mathbb{P}}

% ---------------------------------------------------------------------------
% 定理类环境：与原文一致——所有定理类环境共享一个计数器，按节编号；
% 原文用 \swapnumbers（编号在前），中译本按中文惯例作"定理 1.1"。
% ---------------------------------------------------------------------------
\newtheoremstyle{cnth}%
  {6pt}% above
  {6pt}% below
  {\normalfont}% body
  {}% indent
  {\bfseries}% head font
  {.}% punctuation
  {0.5em}% space after head
  {}% head spec

\theoremstyle{cnth}
\newtheorem{thm}{定理}[section]
\newtheorem{dfn}[thm]{定义}
\newtheorem{defn}[thm]{定义}
\newtheorem{prp}[thm]{命题}
\newtheorem{prop}[thm]{命题}
\newtheorem{crl}[thm]{推论}
\newtheorem{cor}[thm]{推论}
\newtheorem{conj}[thm]{猜想}
\newtheorem{nota}[thm]{记号}
\newtheorem{ques}[thm]{问题}
\newtheorem{exmp}[thm]{例}
\newtheorem{example}[thm]{例}
\newtheorem{lmm}[thm]{引理}
\newtheorem{lemma}[thm]{引理}
\newtheorem{rmk}[thm]{注}
\newtheorem{cnv}[thm]{约定}
\newtheorem{exer}[thm]{练习}

% ---------------------------------------------------------------------------
% 三类译者方框：yizhu（橙色，编号）、conceptbox（紫色，不编号）、
% keytakeaways（蓝色）。三个 \ifshow* 开关可一键关闭：关闭后方框消失、
% 文字以内文排印。
% ---------------------------------------------------------------------------
\newif\ifshowyizhu
\showyizhutrue
\newcounter{yizhu}

\newenvironment{yizhu}[1][]{%
    \refstepcounter{yizhu}%
    \par\medskip
    \ifshowyizhu
        \begin{tcolorbox}[enhanced, breakable,
            colback=orange!4!white, colframe=orange!55!black,
            boxrule=0.5pt, arc=2pt, left=6pt, right=6pt, top=4pt, bottom=4pt,
            title={\small\textbf{译注 \theyizhu\if\relax\detokenize{#1}\relax\else\ ——#1\fi}},
            fonttitle=\bfseries, coltitle=black, colbacktitle=orange!12!white,
            toptitle=1pt, bottomtitle=1pt]
    \fi
}{%
    \ifshowyizhu
        \end{tcolorbox}%
    \fi
    \par\smallskip
}

\newif\ifshowconcept
\showconcepttrue

\newenvironment{conceptbox}[1]{%
    \par\medskip
    \ifshowconcept
        \begin{tcolorbox}[enhanced, breakable,
            colback=violet!3!white, colframe=violet!55!black,
            boxrule=0.5pt, arc=2pt, left=6pt, right=6pt, top=4pt, bottom=4pt,
            title={\small\textbf{概念回顾：#1}\quad{\footnotesize（译者所加）}},
            fonttitle=\bfseries, coltitle=black, colbacktitle=violet!12!white,
            toptitle=1pt, bottomtitle=1pt]
    \fi
}{%
    \ifshowconcept
        \end{tcolorbox}%
    \fi
    \par\medskip
}

\newif\ifshowkey
\showkeytrue

\newenvironment{keytakeaways}{%
    \par\medskip
    \ifshowkey
        \begin{tcolorbox}[enhanced, breakable,
            colback=blue!3!white, colframe=blue!35!black,
            boxrule=0.5pt, arc=2pt, left=6pt, right=6pt, top=4pt, bottom=4pt,
            title={\small\textbf{本节重点}\quad{\footnotesize（译者所加）}},
            fonttitle=\bfseries, coltitle=black, colbacktitle=blue!10!white,
            toptitle=1pt, bottomtitle=1pt]
    \fi
    \begin{itemize}[leftmargin=1.4em,itemsep=0.15em,topsep=2pt]
}{%
    \end{itemize}%
    \ifshowkey
        \end{tcolorbox}%
    \fi
    \par\medskip
}
